% Options for packages loaded elsewhere
\PassOptionsToPackage{unicode}{hyperref}
\PassOptionsToPackage{hyphens}{url}
%
\documentclass[
]{article}
\usepackage{amsmath,amssymb}
\usepackage{iftex}
\ifPDFTeX
  \usepackage[T1]{fontenc}
  \usepackage[utf8]{inputenc}
  \usepackage{textcomp} % provide euro and other symbols
\else % if luatex or xetex
  \usepackage{unicode-math} % this also loads fontspec
  \defaultfontfeatures{Scale=MatchLowercase}
  \defaultfontfeatures[\rmfamily]{Ligatures=TeX,Scale=1}
\fi
\usepackage{lmodern}
\ifPDFTeX\else
  % xetex/luatex font selection
\fi
% Use upquote if available, for straight quotes in verbatim environments
\IfFileExists{upquote.sty}{\usepackage{upquote}}{}
\IfFileExists{microtype.sty}{% use microtype if available
  \usepackage[]{microtype}
  \UseMicrotypeSet[protrusion]{basicmath} % disable protrusion for tt fonts
}{}
\makeatletter
\@ifundefined{KOMAClassName}{% if non-KOMA class
  \IfFileExists{parskip.sty}{%
    \usepackage{parskip}
  }{% else
    \setlength{\parindent}{0pt}
    \setlength{\parskip}{6pt plus 2pt minus 1pt}}
}{% if KOMA class
  \KOMAoptions{parskip=half}}
\makeatother
\usepackage{xcolor}
\setlength{\emergencystretch}{3em} % prevent overfull lines
\providecommand{\tightlist}{%
  \setlength{\itemsep}{0pt}\setlength{\parskip}{0pt}}
\setcounter{secnumdepth}{-\maxdimen} % remove section numbering
\ifLuaTeX
  \usepackage{selnolig}  % disable illegal ligatures
\fi
\IfFileExists{bookmark.sty}{\usepackage{bookmark}}{\usepackage{hyperref}}
\IfFileExists{xurl.sty}{\usepackage{xurl}}{} % add URL line breaks if available
\urlstyle{same}
\hypersetup{
  pdftitle={Lenin 1905rd The Plan of the St. Petersburg Battle},
  pdfauthor={V.I. Lenin},
  hidelinks,
  pdfcreator={LaTeX via pandoc}}

\title{Lenin 1905rd The Plan of the St. Petersburg Battle}
\author{V.I. Lenin}
\date{}

\begin{document}
\maketitle

\hypertarget{lenin-1905rd-the-plan-of-the-st.-petersburg-battle}{%
\section{Lenin: 1905/rd: The Plan of the St. Petersburg
Battle}\label{lenin-1905rd-the-plan-of-the-st.-petersburg-battle}}

\begin{quote}
\hypertarget{excerpt}{%
\subsection{Excerpt}\label{excerpt}}

The Plan of the St. Petersburg Battle
\end{quote}

\begin{center}\rule{0.5\linewidth}{0.5pt}\end{center}

\hypertarget{v.-i.-lenin}{%
\subsection{\texorpdfstring{\href{https://www.marxists.org/archive/lenin/works/1905/jan/25.htm}{V.
I.} ~
\href{https://www.marxists.org/archive/lenin/works/1905/jan/19b.htm}{Lenin}}{V. I. ~ Lenin}}\label{v.-i.-lenin}}

\hypertarget{revolutionary-days}{%
\subsubsection{\texorpdfstring{\href{https://www.marxists.org/archive/lenin/works/1905/rd/index.htm\#3}{Revolutionary
Days}}{Revolutionary Days}}\label{revolutionary-days}}

\begin{center}\rule{0.5\linewidth}{0.5pt}\end{center}

\hypertarget{the-plan-of-the-st.-petersburg-battle}{%
\paragraph{The Plan of the St. Petersburg
Battle}\label{the-plan-of-the-st.-petersburg-battle}}

It seems strange, at first glance, to refer to the peaceful march of
unarmed workers to present a petition as a battle. It was a massacre.
But the government had looked forward to a battle, and it doubtlessly
acted according to a well-laid plan. It considered the defense of St.
Petersburg and of the Winter Palace from the military standpoint. It
took all necessary military measures. It removed all the civil
authorities, and placed the capital with its million and a half
population under the complete control of the generals (headed by Grand
Duke Vladimir), who were thirsting for the blood of the people.

The government deliberately drove the proletariat to revolt, provoked
it, by the massacre of unarmed people, to erect barricades, in order to
drown the uprising in a sea of blood. The proletariat will learn from
these military lessons afforded by the government. For one thing, it
will learn the art of civil war, now that it has started the revolution.
Revolution is war. Of all the wars known in his tory it is the only
lawful, rightful, just, and truly great war. This war is not waged in
the selfish interests of a handful of rulers and exploiters, like any
and all other wars, but in the interests of the masses of the people
against the tyrants, in the interests of the toiling and exploited
millions upon millions against despotism and violence.

All detached observers now are of one accord in admitting that in Russia
this war has been declared and begun. The proletariat will rise again in
still greater masses. What is left of the childish faith in the tsar
will now vanish as quickly as the St. Petersburg workers changed from ~
petitioning to barricade fighting. The workers everywhere will arm. What
matters it that the police will keep a tenfold greater watch over the
arsenals and arms stores and shops? No stringencies, no prohibitions
will stop the masses in the cities, once they have come to realize that
without arms they can always be shot down by the government on the
slightest pretext. Everyone will try his hardest to get him self a gun
or at least a revolver, to conceal his fire-arms from the police and be
ready to repel any attack of the blood thirsty servitors of tsarism.
Every beginning is difficult, as the saying goes. It was very difficult
for the workers to go over to the armed combat. The government has now
forced them to it. The first and most difficult step has been taken.

An English correspondent reports a typical conversation among workers in
a Moscow street. A group of workers was openly discussing the lessons of
the day. ``Hatchets?'' said one. ``No, you can't do anything with a
hatchet against a sabre. You can't get at him with a hatchet any more
than you can with a knife. No, what we need is revolvers, revolvers at
the very least, and better still, guns.'' Such conversations can be
heard now all over Russia. And these conversations after ``Vladimir's
Day'' in St. Petersburg will not remain mere talk.

The military plan of the tsar's uncle, Vladimir, who directed the
massacre, was to keep the people from the suburbs, the workers' suburbs,
away from the centre of the city. No pains were spared to make the
soldiers believe that the workers wanted to demolish the Winter Palace
(by means of icons, crosses, and petitions!) and kill the tsar. The
strategic task was simply to guard the bridges and the main streets
leading to the Palace Square. And the principal scenes of ``military
operations'' were the squares near the bridges (the Troitsky,
Samsonievsky, Nikolayevsky, and Palace bridges), as well as the streets
leading from the working-class districts to the centre (the Narvskaya
Zastava, Schlüsselburg Highway, and Nevsky Prospekt), and, lastly, the
Palace Square itself, to which thousands upon thousands of workers
penetrated in spite of the massed troops and the resistance they met
with. Military operations were, of course, rendered much easier by the
fact that everybody ~ knew perfectly well where the workers were going,
that there was but one rallying point and one. objective. The valiant
generals attacked ``successfully'' an enemy who had come unarmed and
made his destination and purpose known in advance...: It was a
dastardly, cold-blooded massacre of defenceless and peaceful people. For
a long time to come new the masses will think over and re live in memory
and in story all that took place. The sole and inevitable conclusion
drawn from these reflections, from the assimilation of ``Vladimir's
lesson'' in the minds of the masses, will be \emph{à la guerre comme à
la guerre}. The working-class masses, and, following their lead, the
masses of the rural poor, will realise that they are combatants in a
war, and then ... then the next battles of our civil war will be fought
according to plan, but no longer according to the ``plan'' of grand
dukes and the tsars. The call ``To arms!'' which sounded among a crowd
of workers in Nevsky Prospekt on January 9 cannot die away now without
reverberation.

\begin{center}\rule{0.5\linewidth}{0.5pt}\end{center}

~

\hypertarget{notes}{%
\subsubsection{Notes}\label{notes}}

\begin{center}\rule{0.5\linewidth}{0.5pt}\end{center}

\end{document}
